\documentclass[popl.tex]{subfiles}
\begin{document}

\clearpage\section{Related Work}\label{sec:related}

Our work draws a threefold correspondence between well-known techniques in (1) formal language theory, (2) program analysis and (3) incremental decoding. We will first survey these topics, then turn our attention to machine learning, with which we compare and partly use for reranking.

\subsection{Formal language theory}

Context-free language (CFL) parsing is the well-studied problem of how to turn a string into a unique tree, with many different algorithms and implementations (e.g., shift-reduce, recursive-descent, LR). Many of those algorithms expect grammars to be expressed in a certain form (e.g., left- or right-recursive) or are optimized for a narrow class of grammars (e.g., regular, linear).

General CFL parsing allows ambiguity (non-unique trees) and can be formulated as a dynamic programming problem, as shown by Cocke-Younger-Kasami (CYK)~\cite{sakai1961syntax}, Earley~\cite{earley1970efficient} and others. These parsers have roughly cubic complexity with respect to the length of the input string.

As shown by Valiant~\cite{valiant1975general}, Lee~\cite{lee2002fast} and others, general CFL recognition is in some sense equivalent to binary matrix multiplication, another well-studied combinatorial problem, known to be at worst subcubic. This reduction opens the door to a range of complexity-theoretic speedups to CFL recognition; however large constants tend to limit their practical applicability. In the parallel setting, Ruzzo~\cite{ruzzo1978general} presents a polylogarithmic parallel parser for length-\(n\) strings, however the overhead can be substantial. Brent and Goldschlager~\cite{brent1984parallel} independently show an approach requiring $\mathcal{O}(n^6)$ processors, which Mikhelson and Okhotin~\cite{mikhelson2025parallel} later extend to weighted semirings. Instead, we opt for a na\"ive parallel CKY algorithm, which is linear time decidable on $\mathcal{O}(n^3)$ processors.

Bar-Hillel~\cite{bar1961formal} proves the closure of CFLs under intersection with regular languages, but does not elaborate on how to construct the corresponding grammar. Salomaa~\cite{salomaa1973formal} and Pasti et al.~\cite{pasti2023intersection} provide helpful insights into constructing the intersection grammar, and Nederhof and Satta~\cite{nederhof2004language} specifically consider finite CFL intersections, but seem unaware of the connection to CFL reachability. Our work specializes Bar-Hillel intersections to Levenshtein automata in particular, and more generally acyclic automata using a refinement of Salomaa's construction~\cite{salomaa1973formal} based on CFL reachability.

\subsection{CFL reachability}

Our contribution is closely related to the literature on CFL reachability. In brief, the CFL reachability problem seeks to determine, given an edge-labeled graph and distinguished vertex pair, $\langle v, v' \rangle$, whether there is a path, $v \rightsquigarrow v'$, whose concatenated edge labels are contained in the CFL. For a deeper overview, see Zhang and Su~\cite{zhang2017context}. This problem has been known~\cite{reps1998program} for some time~\cite{kodumal2004set} to have broad applications to program analysis and as our work finds, to syntactic program repair.

From a complexity-theoretic perspective, the CFL reachability problem is known to be at worst subcubic~\cite{chistikov2022subcubic} with polylogarithmic time factors. Koutris and Deep~\cite{koutris2023fine} present a fine-grained complexity analysis, with concurrent work by Istomina et al.~\cite{istomina2023fine} expanding on fine-grained reductions. Muravev and Grigorev explore how to accelerate this technique on a GPU~\cite{muravev2025universal}.

Surprisingly absent from the literature on CFL reachability is a discussion of the Bar-Hillel construction, regular expressions for witnessability, or the use of Brzozowski's derivative for incremental decoding. Nor does the literature specifically consider finite intersection nonemptiness between CFLs and acyclic automata such as the Levenshtein NFA (\S~\ref{sec:repair_ex}). In Theorem~\ref{thm:parallel_decision_complexity} we give a constructive proof via a parallel matrix closure with a linear bound on the number of squaring rounds, borrowing the matrix multiplication technique from CFL reachability to build a star-free regular expression that we decode using the Brzozowski derivative. This technique sheds new light on the Bar-Hillel construction, and translates to a simple and efficient implementation which is fully compatible with autoregressive decoding techniques used in probabilistic language modeling.

\subsection{Language equations}

Language equations are a powerful tool for reasoning about formal languages and their inhabitants. First proposed by Ginsburg et al.~\cite{ginsburg1962two} for the ALGOL language, language equations are essentially systems of inequalities with variables representing \textit{holes}, i.e., unknown values, in the language or grammar. Solutions to these equations can be obtained using various fixpoint techniques, yielding members of the language. This insight reveals the true algebraic nature of CFLs and their cousins.

Being an algebraic formalism, language equations naturally give rise to a kind of calculus, vaguely reminiscent of Leibniz's and Newton's. First studied by Brzozowski~\cite{brzozowski1964derivatives, brzozowski1980equations} and Antimirov~\cite{antimirov1996partial}, one can take the derivative of a language equation, which can be interpreted as a kind of continuation or language quotient, revealing the suffixes that complete a given prefix. This technique leads to an elegant family of algorithms for incremental parsing~\cite{might2011parsing, adams2016complexity} and regular expression matching~\cite{stanford2021symbolic,varatalu2025re}.

Recently Pasti et al.~\cite{pasti2026prefix} introduce an elegant grammar-to-grammar transformation and next-token decoding technique that gives the exact next-token probabilities for a given PCFG, across all prefixes. This result is a surprising one, because the weighted algorithm was only recently discovered, despite PCFGs and autoregressive decoding techniques being used for over 70 years.

%From a more applied perspective, parsers are ubiquitous in software engineering, but none are designed to handle arbitrary CFGs or recover from arbitrary errors. Parr and Quong introduce ANTLR~\cite{parr1995antlr} which can handle LL(k) grammars and offers an IDE plugin with limited support for error recovery. Scott and Johnstone~\cite{scott2010gll} introduce GLL parsing, which supports linear-time parsing for LL grammars and cubic for arbitrary CFGs, but has limited support for error recovery. Inspired by their work, we introduce a method for repairing small syntax errors in arbitrary CFLs.

More concretely, we restrict our attention to language equations over CFLs whose variables coincide with edit locations in the source code of a computer program, and solutions correspond to syntax repairs. While prior work has studied the use of language equations for parsing~\cite{might2011parsing}, to our knowledge, they were never specifically considered for code completion or syntax error correction.

\subsection{Syntax repair}

In finite languages, syntax repair corresponds to spelling correction, a more restrictive and largely solved problem. Schulz and Mihov~\cite{schulz2002fast} construct a finite automaton that returns the nearest dictionary entry by Levenshtein edit distance. Though considerably simpler than syntax correction, their work shares similar challenges and offers insights for handling more general repair scenarios.

When a string is syntactically invalid, parsing grows more challenging. Like spelling, the problem is to find the minimum number of edits required to transform an arbitrary string into a syntactically valid one, where validity is defined as containment in a (typically) context-free language. Early work, including Irons~\cite{irons1963error} and Aho~\cite{aho1972minimum} proposes a dynamic programming algorithm to compute the minimum number of edits required to fix an invalid string. Prior work on error-correcting parsers only considers the nearest edit(s), and does not study edits of varying distance in the Levenshtein ball. Furthermore, the problem of repair is not generally well-posed, as there can be many valid solutions. We instead focus on maximum probability Levenshtein-CFL reachability, which seeks to find the most natural syntactically valid repair within a fixed Levenshtein distance.

Diekmann and Tratt~\cite{diekmann2018dont} present a rule-based syntax repair tool that retrieves the complete set of minimum-cost repairs, but only works for deterministic CFLs, a proper subset of the CFL family which admit a linear-time parser. Their cost model is based on insertion and deletion, and does not consider probability or non-minimal edit distance. Tidyparse can handle arbitrary CFLs and generate repairs within an arbitrary edit distance, using a Levenshtein cost model.

Zhang et al.~\cite{zhang2023ordinalfix} introduce OrdinalFix, which uses CFL reachability to repair compiler errors, however their method only returns admissible repairs and not necessarily probable ones. As they do not consider the problem of maximum-probability repairs, nor use any form of ranking to sort the results by naturalness or probability, we do not compare with their work.

%\subsection{String solving}
%
%There is related work on string constraints in the constraint programming literature, featuring solvers like CFGAnalyzer and HAMPI~\cite{kiezun2009hampi}, which consider bounded context free grammars and intersections thereof. Boja{\'n}czyk et al. (2014)~\cite{bojanczyk2014automata} introduce the theory of nominal automata. Around the same time, D'Antoni et al. (2014) introduce \textit{symbolic automata}~\cite{dantoni2014minimization}, a generalization of finite automata which allow infinite alphabets and symbolic expressions over them. Hague et al. (2024)~\cite{hague2024parikh} use Parikh's theorem in the context of symbolic automata to speed up string constraint solving. In none of the constraint programming literature we surveyed do any of the approaches specifically consider the problem of syntax error correction, which is the core focus of our work.

%\subsection{Error correcting codes}
%
%Our work focuses on errors arising from human factors in computer programming, in particular \textit{syntax error correction}, which is the problem of fixing partially corrupted programs. Modern research on error correction, however, can be traced back to the early days of coding theory when researchers designed \textit{error-correcting codes} (ECCs) to denoise transmission errors induced by external interference, e.g., collision with a high-energy proton, manipulation by an adversary or even typographical mistake. In this context, \textit{code} can be any logical representation for communicating information between two parties (such as a human and a computer), and an ECC is a carefully-designed scheme which ensures that even if some portion of the message should become corrupted, one can still recover the original message by solving a linear system of equations. When designing ECCs, one typically assumes a noise model over a certain sample space, such as the Hamming~\cite{titsias2017hamming} or Levenshtein~\cite{levenshtein1966binary, becerra2008learning, barlev2021levenshtein} balls, from which we draw inspiration for this work.

\subsection{Decoding}\label{sec:decoding}

Decoding is an important subproblem in machine translation, speech recognition, and sequence-to-sequence learning. Given a compressed encoding of some learned distribution, the goal is to find the maximum probability samples. A classic example is Viterbi decoding, used to find the most likely path through a hidden Markov model (HMM), a kind of weighted automaton.

In particular, we care about the problem of \textit{top-$k$ decoding}, which attempts to find the exact or approximate $k$-most likely samples in order of decreasing probability. This is closely related to the $k$-best enumeration~\cite{eppstein2014k} problem, a carefully studied problem in graph theory and combinatorial optimization. An exact solution to this problem for large acyclic CFGs is often intractable, but we can approximate it using a beam search over subwords, then rerank top-scoring results.

A popular solution to $k$-best decoding in the NLP literature is a technique called cube-pruning~\cite{huang2005better, chiang2007hierarchical}, which samples maximum probability paths through a hypergraph. We take inspiration from this technique and adapt it to the setting of constrained decoding from finite CFGs.

An alternate line of work originates from combinatorics~\cite{hickey1983uniform} and Boltzmann sampling~\cite{duchon2004boltzmann}, which constructs a generating function for the language and samples it uniformly. This technique has applications to constraint satisfaction and model counting problems in formal languages.

A third approach would be to use some form of constrained decoding~\cite{willard2023efficient, ugare2024improving, loula2025syntactic} such as sequential Monte Carlo to steer an autoregressive LLM, as proposed by Lew et al.~\cite{lew2023sequential, lipkin2025fast}. These techniques show promise for program repair, however, the question of whether to use left-to-right decoding or some other strategy is still unresolved in the language modeling community. For example, there is an emerging class of flow-based or structured denoising diffusion models~\cite{austin2021structured} which starts from a noise distribution and iteratively denoises it by sampling one or more edits at random locations with each decoding step. Typical work focuses on audiovisual data, but very recent work by Havasi et al.~\cite{havasi2025edit} adapt this to the Levenshtein edit model for generating source code. Although these models do not yet use CFGs or consider language intersections, they are inherently more fault-tolerant than decoders which require expensive backtracking-style search.

\subsection{Learning-based program repair methods}

The last decade has seen a surge of progress in programming with large language models. That work is primarily based on methods from differential calculus and continuous optimization, leading to the so-called \textit{naturalness hypothesis}~\cite{allamanis2018survey}, which suggests programming languages are not so different from natural ones. In contrast, PL theory takes the view that languages are essentially discrete sets governed by logical calculi. Programming, thus viewed, is more like a mathematical exercise in constraint satisfaction. These two approaches have more in common than would seem.

From an applied perspective, a number of gradient-based methods have been introduced to repair programming errors~\cite{allamanis2021self, chirkova2021empirical, drain2021generating}. These approaches typically employ large language models (LLMs) and treat the problem as a sequence-to-sequence transformation. While capable of generating natural repairs, these models are susceptible to misgeneralization, costly to train, and challenging to customize thereafter. Furthermore, the generated repairs are not necessarily sound without additional filtering, and we observe the released models often hallucinate false-positive repairs.

Two prior works specifically address syntax repair, Break-It-Fix-It (BIFI)~\cite{yasunaga2021break} and Seq2Parse~\cite{sakkas2022seq2parse}. BIFI adapts techniques from semi-supervised learning to generate synthetic errors in clean code and fixes them. This reduces the need for pairwise training data, but generalizes poorly to lengthy or out-of-distribution repairs. Seq2Parse combines a transformer-based model with an augmented version of the Earley parser to suggest error rules, but only suggests a single repair.

Recent work by Merrill et al.~\cite{merrill2022saturated} and Chiang et al.~\cite{chiang2023tighter} suggests that the issue with generalization may be more foundational: transformer-based language models, a popular class of neural language models used in program synthesis and repair, are fundamentally less expressive than context-free grammars, which formally describe the syntax of many programming languages. This suggests such models, despite their useful approximation properties, are ill-suited for the task of end-to-end syntax repair. Yet, as our work demonstrates, they can be useful for resolving ambiguity between valid repairs of differing probability or reranking a set of repair candidates drawn from a CFL.

Using the RASP model from Weiss et al.~\cite{weiss2021thinking}, Yang et al.~\cite{yang2024masked} characterize the expressive power of hard-attention transformers in terms of star-free regular expressions. While their work uses formal language theory to investigate language learnability and structural priors in transformers, it shares fruitful connections to CFL reachability, which, similar to RASP, treats matrix multiplication as a kind of programmable interface into which various state tracking problems and static analysis tasks can be compiled and analyzed in a common linear algebraic framework.
\ifSubfilesClassLoaded{\bibliography{../bib/acmart}}{}
\end{document}
